% =============================================================================
% Chapter 10 — The Expansion Memory System ($BA00–$BFFF)
% =============================================================================
\chapter{The Expansion Memory System}

\epigraph{``640K ought to be enough for anybody --- until you actually try to
do something with it.''}{\textit{--- Apocryphal, circa 1981}}

\section*{Introduction}

The 96-byte address range from \$BA00 to \$BFFF packs five distinct
subsystems into a dense block of hardware registers, followed by a 1\,KB
memory-mapped window into expansion RAM:

\begin{enumerate}
  \item \textbf{Expansion Memory Controller} (\$BA00--\$BA3F) --- manages
        512\,KB of banked XRAM with DMA transfers and a named-block filesystem.
  \item \textbf{Timer Controller} (\$BA40--\$BA4F) --- a programmable interval
        timer that can generate IRQs at configurable rates.
  \item \textbf{Music Status Registers} (\$BA50--\$BA62) --- read-only
        registers exposing the state of the 14-voice music engine.
  \item \textbf{DMA Controller} (\$BA63--\$BA82) --- a high-speed block
        transfer engine that moves data between any combination of CPU RAM,
        VGC memory spaces, and XRAM.
  \item \textbf{Blitter Controller} (\$BA83--\$BAA2) --- a 2D block transfer
        engine with stride control, fill mode, and color-key transparency.
\end{enumerate}

The final 1\,KB at \$BC00--\$BFFF provides four 256-byte windows that
map directly into XRAM, allowing the 6502 to read and write expansion
memory as though it were ordinary RAM.


% ═══════════════════════════════════════════════════════════════════════════════
\section{Expansion Memory Controller (\$BA00--\$BA3F)}
\label{sec:xmc-registers}
% ═══════════════════════════════════════════════════════════════════════════════

The Expansion Memory Controller (XMC) provides access to 512\,KB of
expansion RAM (XRAM) organized as 8 banks of 64\,KB each.  Data moves
between CPU RAM and XRAM via Stash/Fetch commands, or byte-at-a-time
through the data port.  A named-block subsystem provides a simple filesystem
within XRAM, and page-level allocation tracking prevents accidental
overwrites.

\subsection{Register Map}

\begin{center}
\small
\begin{longtable}{@{}llcll@{}}
\toprule
\textbf{Address} & \textbf{Name} & \textbf{R/W} & \textbf{Default} & \textbf{Description} \\
\midrule
\endfirsthead
\toprule
\textbf{Address} & \textbf{Name} & \textbf{R/W} & \textbf{Default} & \textbf{Description} \\
\midrule
\endhead
\$BA00 & XmcCmd       & W   & ---  & Command register \\
\$BA01 & XmcStatus    & R   & \$00 & Status code \\
\$BA02 & XmcErrCode   & R   & 0    & Error code \\
\$BA03 & XmcCfg       & R/W & 0    & Reserved \\
\$BA04 & XmcAddrL     & R/W & 0    & XRAM address low byte \\
\$BA05 & XmcAddrM     & R/W & 0    & XRAM address middle byte \\
\$BA06 & XmcAddrH     & R/W & 0    & XRAM address high byte \\
\$BA07 & XmcRamL      & R/W & 0    & CPU RAM address low byte \\
\$BA08 & XmcRamH      & R/W & 0    & CPU RAM address high byte \\
\$BA09 & XmcLenL      & R/W & 0    & Transfer length low byte \\
\$BA0A & XmcLenH      & R/W & 0    & Transfer length high byte \\
\$BA0B & XmcData      & R/W & 0    & Single-byte data port \\
\$BA0C & XmcBank      & R/W & 0    & Active 64\,KB bank selector \\
\$BA0D & XmcBanks     & R   & 8    & Total bank count \\
\$BA0E & XmcPagesUsedL & R  & 0    & Used 256-byte pages (low) \\
\$BA0F & XmcPagesUsedH & R  & 0    & Used 256-byte pages (high) \\
\$BA10 & XmcPagesFreeL & R  & 0    & Free pages (low) \\
\$BA11 & XmcPagesFreeH & R  & 0    & Free pages (high) \\
\$BA12 & XmcNameLen   & R/W & 0    & Named block name length (0--27) \\
\$BA13 & XmcHandle    & R/W & 0    & Named block handle \\
\$BA14 & XmcDirCountL & R   & 0    & Named block count (low) \\
\$BA15 & XmcDirCountH & R   & 0    & Named block count (high) \\
\$BA16 & XmcWinCtl    & R/W & 0    & Window enable bitmask (bits 0--3) \\
\$BA18--\$BA1A & XmcWin0A & R/W & 0 & Window 0 XRAM base (24-bit LE) \\
\$BA1B--\$BA1D & XmcWin1A & R/W & 0 & Window 1 XRAM base (24-bit LE) \\
\$BA1E--\$BA20 & XmcWin2A & R/W & 0 & Window 2 XRAM base (24-bit LE) \\
\$BA21--\$BA23 & XmcWin3A & R/W & 0 & Window 3 XRAM base (24-bit LE) \\
\$BA24--\$BA3F & XmcName  & R/W & 0 & ASCII name buffer (28 bytes) \\
\bottomrule
\end{longtable}
\end{center}


\subsection{Detailed Register Descriptions}

% --- $BA00 XmcCmd --------------------------------------------------------

\mementry{BA00}{XmcCmd}{W}{---}

Command register.  Writing a command byte here immediately executes the
corresponding XMC operation.  All parameters (addresses, lengths, names)
must be set in the surrounding registers \emph{before} writing to XmcCmd.
See Section~\ref{sec:xmc-commands} for the full command reference.

% --- $BA01 XmcStatus -----------------------------------------------------

\mementry{BA01}{XmcStatus}{R}{\$00}

Status of the last command:

\begin{center}
\small
\begin{tabular}{@{}cl@{}}
\toprule
\textbf{Value} & \textbf{Meaning} \\
\midrule
\$00 & Idle --- no command has been issued \\
\$02 & OK --- last command completed successfully \\
\$03 & Error --- last command failed (see XmcErrCode) \\
\bottomrule
\end{tabular}
\end{center}

% --- $BA02 XmcErrCode ----------------------------------------------------

\mementry{BA02}{XmcErrCode}{R}{0}

Detailed error code when XmcStatus is \$03:

\begin{center}
\small
\begin{tabular}{@{}cl@{}}
\toprule
\textbf{Value} & \textbf{Meaning} \\
\midrule
0 & No error \\
1 & Range error --- address or length out of bounds \\
2 & Bad arguments --- invalid parameter combination \\
3 & Not found --- named block does not exist \\
4 & No space --- insufficient free XRAM for allocation \\
5 & Bad name --- invalid or empty name string \\
6 & End of directory --- no more named blocks to enumerate \\
\bottomrule
\end{tabular}
\end{center}

% --- $BA03 XmcCfg --------------------------------------------------------

\mementry{BA03}{XmcCfg}{R/W}{0}

Reserved for future use.  Reads return 0; writes are accepted but have no
effect.

% --- $BA04-$BA06 XmcAddr -------------------------------------------------

\mementry{BA04}{XmcAddrL}{R/W}{0}

Low byte of the 24-bit XRAM address.  Together with XmcAddrM (\addr{BA05})
and XmcAddrH (\addr{BA06}), this forms the little-endian target address
within expansion memory.  The full range is \$000000--\$07FFFF (512\,KB).

\mementry{BA05}{XmcAddrM}{R/W}{0}

Middle byte of the 24-bit XRAM address.

\mementry{BA06}{XmcAddrH}{R/W}{0}

High byte of the 24-bit XRAM address.

\begin{tipbox}
To set the XRAM address to \$012345, write \$45 to \addr{BA04}, \$23 to
\addr{BA05}, and \$01 to \addr{BA06}.  Little-endian byte order throughout.
\end{tipbox}

% --- $BA07-$BA08 XmcRam --------------------------------------------------

\mementry{BA07}{XmcRamL}{R/W}{0}

Low byte of the 16-bit CPU RAM address used by Stash and Fetch commands.

\mementry{BA08}{XmcRamH}{R/W}{0}

High byte of the CPU RAM address.

% --- $BA09-$BA0A XmcLen --------------------------------------------------

\mementry{BA09}{XmcLenL}{R/W}{0}

Low byte of the 16-bit transfer length in bytes.

\mementry{BA0A}{XmcLenH}{R/W}{0}

High byte of the transfer length.

% --- $BA0B XmcData -------------------------------------------------------

\mementry{BA0B}{XmcData}{R/W}{0}

Single-byte data port.  Used by the GetByte and PutByte commands to transfer
one byte at a time between the 6502 and XRAM at the address specified by
XmcAddr.  Also used by the Fill command as the fill value.

% --- $BA0C XmcBank -------------------------------------------------------

\mementry{BA0C}{XmcBank}{R/W}{0}

Active 64\,KB bank selector.  Valid range: 0--7 (for the standard 512\,KB
configuration).  This register provides a convenient shorthand for setting
the high byte of XmcAddr: writing bank $n$ is equivalent to setting
XmcAddrH to $n$ and clearing XmcAddrM and XmcAddrL.

\begin{notebox}
The bank selector is a convenience for BASIC programs that think in terms
of 64\,KB banks.  Assembly programs typically address XRAM directly via the
24-bit XmcAddr registers.
\end{notebox}

% --- $BA0D XmcBanks ------------------------------------------------------

\mementry{BA0D}{XmcBanks}{R}{8}

Total number of 64\,KB banks available.  Typically 8, giving 512\,KB of
expansion memory.  This is a read-only hardware configuration value.

% --- $BA0E-$BA0F XmcPagesUsed -------------------------------------------

\mementry{BA0E}{XmcPagesUsedL}{R}{0}

Low byte of the number of 256-byte pages currently marked as used.

\mementry{BA0F}{XmcPagesUsedH}{R}{0}

High byte of the used page count.  Issue the Stats command (\$07) to
refresh these values before reading.

% --- $BA10-$BA11 XmcPagesFree -------------------------------------------

\mementry{BA10}{XmcPagesFreeL}{R}{0}

Low byte of the number of 256-byte pages currently free.

\mementry{BA11}{XmcPagesFreeH}{R}{0}

High byte of the free page count.

% --- $BA12 XmcNameLen ---------------------------------------------------

\mementry{BA12}{XmcNameLen}{R/W}{0}

Length of the name string in the XmcName buffer.  Valid range: 0--27.  Set
this before issuing named-block commands (NStash, NFetch, NDelete,
NDirOpen, NDirRead).

% --- $BA13 XmcHandle ----------------------------------------------------

\mementry{BA13}{XmcHandle}{R/W}{0}

Handle for the current named block.  After a successful NStash or NFetch,
this register holds the handle assigned to (or looked up for) the block.
The handle can be used for subsequent operations without re-specifying the
name.

% --- $BA14-$BA15 XmcDirCount -------------------------------------------

\mementry{BA14}{XmcDirCountL}{R}{0}

Low byte of the total number of named blocks currently stored in XRAM.

\mementry{BA15}{XmcDirCountH}{R}{0}

High byte of the named block count.

% --- $BA16 XmcWinCtl ---------------------------------------------------

\mementry{BA16}{XmcWinCtl}{R/W}{0}

Window enable bitmask.  Each of the four bits enables one of the 256-byte
memory-mapped windows at \$BC00--\$BFFF:

\bitfield{---}{---}{---}{---}{Win3}{Win2}{Win1}{Win0}

\begin{center}
\small
\begin{tabular}{@{}cl@{}}
\toprule
\textbf{Bit} & \textbf{Window} \\
\midrule
0 & Window 0 (\$BC00--\$BCFF) \\
1 & Window 1 (\$BD00--\$BDFF) \\
2 & Window 2 (\$BE00--\$BEFF) \\
3 & Window 3 (\$BF00--\$BFFF) \\
\bottomrule
\end{tabular}
\end{center}

When a window is disabled, reads return \$00 and writes are ignored.

% --- $BA18-$BA1A XmcWin0A -----------------------------------------------

\memrange{BA18}{BA1A}{XmcWin0A}{R/W}{0}

24-bit little-endian XRAM base address for Window~0.  Reads and writes to
\$BC00--\$BCFF are redirected to XRAM starting at this address.

% --- $BA1B-$BA1D XmcWin1A -----------------------------------------------

\memrange{BA1B}{BA1D}{XmcWin1A}{R/W}{0}

24-bit XRAM base address for Window~1 (\$BD00--\$BDFF).

% --- $BA1E-$BA20 XmcWin2A -----------------------------------------------

\memrange{BA1E}{BA20}{XmcWin2A}{R/W}{0}

24-bit XRAM base address for Window~2 (\$BE00--\$BEFF).

% --- $BA21-$BA23 XmcWin3A -----------------------------------------------

\memrange{BA21}{BA23}{XmcWin3A}{R/W}{0}

24-bit XRAM base address for Window~3 (\$BF00--\$BFFF).

% --- $BA24-$BA3F XmcName ------------------------------------------------

\memrange{BA24}{BA3F}{XmcName}{R/W}{0}

28-byte ASCII name buffer for named-block operations.  Store the block name
here before issuing NStash, NFetch, NDelete, or NDirOpen.  After NDirRead,
this buffer contains the name of the current directory entry.  Names are
null-padded; set XmcNameLen (\addr{BA12}) to the actual string length.


% ─────────────────────────────────────────────────────────────────────────────
\subsection{Command Reference}
\label{sec:xmc-commands}
% ─────────────────────────────────────────────────────────────────────────────

Commands are executed by writing the command byte to \addr{BA00}.  Set all
required parameters \emph{before} writing the command byte.

\subsubsection{Byte-Level Access}

\cmdentry{01}{GetByte}{Read one byte from XRAM}

Reads the byte at the XRAM address specified by XmcAddr
(\addr{BA04}--\addr{BA06}) and places it in XmcData (\addr{BA0B}).

\textbf{Parameters:}
\begin{itemize}
  \item XmcAddr (\$BA04--\$BA06): 24-bit XRAM source address
\end{itemize}

\textbf{Result:} XmcData contains the byte read.  XmcStatus $=$ \$02.

\medskip

\cmdentry{02}{PutByte}{Write one byte to XRAM}

Writes the value in XmcData (\addr{BA0B}) to the XRAM address specified
by XmcAddr (\addr{BA04}--\addr{BA06}).

\textbf{Parameters:}
\begin{itemize}
  \item XmcAddr (\$BA04--\$BA06): 24-bit XRAM destination address
  \item XmcData (\$BA0B): byte to write
\end{itemize}

\textbf{Result:} XmcStatus $=$ \$02.

\subsubsection{Block Transfers}

\cmdentry{03}{Stash}{Copy CPU RAM to XRAM}

Copies XmcLen bytes from CPU RAM starting at XmcRam to XRAM starting at
XmcAddr.

\textbf{Parameters:}
\begin{itemize}
  \item XmcRam (\$BA07--\$BA08): 16-bit source address in CPU RAM
  \item XmcAddr (\$BA04--\$BA06): 24-bit destination address in XRAM
  \item XmcLen (\$BA09--\$BA0A): number of bytes to copy
\end{itemize}

\textbf{Result:} XmcStatus $=$ \$02 on success, \$03 on range error.

\medskip

\cmdentry{04}{Fetch}{Copy XRAM to CPU RAM}

Copies XmcLen bytes from XRAM starting at XmcAddr to CPU RAM starting at
XmcRam.

\textbf{Parameters:}
\begin{itemize}
  \item XmcAddr (\$BA04--\$BA06): 24-bit source address in XRAM
  \item XmcRam (\$BA07--\$BA08): 16-bit destination address in CPU RAM
  \item XmcLen (\$BA09--\$BA0A): number of bytes to copy
\end{itemize}

\textbf{Result:} XmcStatus $=$ \$02 on success, \$03 on range error.

\medskip

\cmdentry{05}{Fill}{Fill XRAM with a byte value}

Fills XmcLen bytes of XRAM starting at XmcAddr with the value in XmcData.

\textbf{Parameters:}
\begin{itemize}
  \item XmcAddr (\$BA04--\$BA06): 24-bit start address in XRAM
  \item XmcLen (\$BA09--\$BA0A): number of bytes to fill
  \item XmcData (\$BA0B): fill value
\end{itemize}

\textbf{Result:} XmcStatus $=$ \$02 on success.

\begin{lstlisting}
10 REM fill the first 256 bytes of XRAM with $AA
20 POKE $BA04, 0 : POKE $BA05, 0 : POKE $BA06, 0
30 DOKE $BA09, 256    : REM length = 256
40 POKE $BA0B, $AA    : REM fill value
50 POKE $BA00, $05    : REM execute Fill
60 IF PEEK($BA01)=$02 THEN PRINT "FILLED OK"
\end{lstlisting}

\subsubsection{Page Management}

\cmdentry{07}{Stats}{Refresh page usage statistics}

Recalculates the used and free page counts and updates XmcPagesUsed
(\addr{BA0E}--\addr{BA0F}) and XmcPagesFree (\addr{BA10}--\addr{BA11}).

\textbf{Parameters:} None.

\textbf{Result:} Page counts updated.  XmcStatus $=$ \$02.

\medskip

\cmdentry{08}{ResetUsage}{Clear page tracking}

Clears all page allocation tracking, marking all XRAM pages as free.  Does
not erase the actual data --- only the bookkeeping is reset.

\textbf{Parameters:} None.

\textbf{Result:} XmcStatus $=$ \$02.

\begin{warningbox}
ResetUsage does not zero out XRAM contents.  Named blocks stored prior to
the reset will still be physically present, but the allocator will treat
their pages as free and may overwrite them.
\end{warningbox}

\medskip

\cmdentry{09}{Release}{Mark XRAM range as free}

Marks the pages covered by the range starting at XmcAddr with length
XmcLen as free.

\textbf{Parameters:}
\begin{itemize}
  \item XmcAddr (\$BA04--\$BA06): start of range to release
  \item XmcLen (\$BA09--\$BA0A): number of bytes to release
\end{itemize}

\textbf{Result:} XmcStatus $=$ \$02.

\medskip

\cmdentry{0A}{Alloc}{Allocate XRAM pages}

Allocates XmcLen bytes of contiguous XRAM.  On success, the starting
address of the allocated block is returned in XmcAddr
(\addr{BA04}--\addr{BA06}).

\textbf{Parameters:}
\begin{itemize}
  \item XmcLen (\$BA09--\$BA0A): number of bytes to allocate
\end{itemize}

\textbf{Result:} XmcAddr updated with the allocated address.  XmcStatus
$=$ \$02 on success, \$03 with ErrCode $=$ 4 (no space) if allocation fails.

\begin{lstlisting}
10 REM allocate 1024 bytes of XRAM
20 DOKE $BA09, 1024   : REM length
30 POKE $BA00, $0A    : REM Alloc
40 IF PEEK($BA01)=$03 THEN PRINT "OUT OF MEMORY" : END
50 A = PEEK($BA04) + PEEK($BA05)*256
60 PRINT "ALLOCATED AT: $"; HEX$(A)
\end{lstlisting}

\subsubsection{Named Blocks}

Named blocks provide a simple key-value filesystem within XRAM.  Each block
is identified by a string name (up to 27 characters) and stores an
arbitrary-length chunk of data.  This allows programs to save and retrieve
data structures by name without managing raw XRAM addresses.

\cmdentry{0B}{NStash}{Store CPU RAM as named block}

Copies XmcLen bytes from CPU RAM at XmcRam into XRAM as a named block.
The name is read from XmcName (\addr{BA24}--\addr{BA3F}) with the length
in XmcNameLen (\addr{BA12}).  If a block with the same name already exists,
it is replaced.

\textbf{Parameters:}
\begin{itemize}
  \item XmcRam (\$BA07--\$BA08): source address in CPU RAM
  \item XmcLen (\$BA09--\$BA0A): number of bytes to store
  \item XmcName (\$BA24+): block name (ASCII)
  \item XmcNameLen (\$BA12): name length
\end{itemize}

\textbf{Result:} XmcHandle (\addr{BA13}) set.  XmcStatus $=$ \$02.

\medskip

\cmdentry{0C}{NFetch}{Retrieve named block to CPU RAM}

Looks up a named block by name and copies its contents to CPU RAM at
XmcRam.

\textbf{Parameters:}
\begin{itemize}
  \item XmcRam (\$BA07--\$BA08): destination address in CPU RAM
  \item XmcName (\$BA24+): block name (ASCII)
  \item XmcNameLen (\$BA12): name length
\end{itemize}

\textbf{Result:} XmcLen updated with the block size.  XmcStatus $=$ \$02
on success, \$03 with ErrCode $=$ 3 (not found) if the name does not exist.

\medskip

\cmdentry{0D}{NDelete}{Delete named block}

Deletes the named block identified by XmcName and frees its XRAM pages.

\textbf{Parameters:}
\begin{itemize}
  \item XmcName (\$BA24+): block name (ASCII)
  \item XmcNameLen (\$BA12): name length
\end{itemize}

\textbf{Result:} XmcStatus $=$ \$02 on success, \$03 with ErrCode $=$ 3
if the name does not exist.

\medskip

\cmdentry{0E}{NDirOpen}{Open named block directory}

Opens the named-block directory for sequential enumeration.  Resets the
internal directory index to the first entry.

\textbf{Parameters:} None.

\textbf{Result:} XmcDirCount (\addr{BA14}--\addr{BA15}) updated with the
total number of named blocks.  XmcStatus $=$ \$02.

\medskip

\cmdentry{0F}{NDirRead}{Read next directory entry}

Reads the next named block from the directory.  The block name is placed
in XmcName, the name length in XmcNameLen, and the block size in XmcLen.
When all entries have been read, XmcStatus is set to \$03 with ErrCode
$=$ 6 (end of directory).

\textbf{Parameters:} None (directory must be opened with NDirOpen first).

\textbf{Result:} XmcName, XmcNameLen, and XmcLen populated, or ErrCode
$=$ 6 at end.

\begin{lstlisting}
10 REM list all named blocks in XRAM
20 POKE $BA00, $0E    : REM NDirOpen
30 N = DEEK($BA14) : PRINT N; " BLOCKS FOUND"
40 FOR I=1 TO N
50   POKE $BA00, $0F  : REM NDirRead
60   IF PEEK($BA01)=$03 THEN I=N : GOTO 120
70   L = PEEK($BA12) : N$ = ""
80   FOR J=0 TO L-1
90     N$ = N$ + CHR$(PEEK($BA24+J))
100  NEXT J
110  PRINT N$; " ("; DEEK($BA09); " BYTES)"
120 NEXT I
\end{lstlisting}


% ═══════════════════════════════════════════════════════════════════════════════
\section{Timer Controller (\$BA40--\$BA4F)}
\label{sec:timer}
% ═══════════════════════════════════════════════════════════════════════════════

The Timer Controller provides a single programmable interval timer that can
generate IRQ interrupts at a configurable rate.  The timer is clocked at
the CPU clock rate divided into 100-cycle quanta.  When the internal counter
reaches the programmed divisor, it sets the IRQ pending flag and resets.

% --- $BA40 TimerCtrl ----------------------------------------------------

\mementry{BA40}{TimerCtrl}{R/W}{0}

Timer control register:

\bitfield{---}{---}{---}{---}{---}{---}{---}{En}

\begin{center}
\small
\begin{tabular}{@{}cl@{}}
\toprule
\textbf{Bit} & \textbf{Function} \\
\midrule
0 & Enable: 1 $=$ timer running, 0 $=$ timer stopped \\
1--7 & Reserved (write 0) \\
\bottomrule
\end{tabular}
\end{center}

% --- $BA41 TimerStatus --------------------------------------------------

\mementry{BA41}{TimerStatus}{R}{0}

Timer status register:

\bitfield{---}{---}{---}{---}{---}{---}{---}{IRQ}

\begin{center}
\small
\begin{tabular}{@{}cl@{}}
\toprule
\textbf{Bit} & \textbf{Function} \\
\midrule
0 & IRQ pending: 1 $=$ timer has fired.  Reading this register clears
    the flag. \\
1--7 & Reserved \\
\bottomrule
\end{tabular}
\end{center}

\begin{notebox}
Reading TimerStatus clears the IRQ pending bit.  In an interrupt handler,
always read this register to acknowledge the timer interrupt and prevent
it from firing again immediately.
\end{notebox}

% --- $BA42-$BA43 TimerDiv -----------------------------------------------

\mementry{BA42}{TimerDivL}{R/W}{0}

Low byte of the 16-bit timer divisor.

\mementry{BA43}{TimerDivH}{R/W}{0}

High byte of the timer divisor.  The timer fires every
\texttt{divisor $\times$ 100} CPU cycles.  For example, a divisor of 100
fires every 10,000 cycles.

\begin{lstlisting}
10 REM set up a timer IRQ at ~60 Hz
20 REM (assumes 1 MHz clock: 1000000/60 = 16667 cycles)
30 REM divisor = 16667/100 = 167
40 DOKE $BA42, 167    : REM set divisor
50 POKE $BA40, 1      : REM enable timer
60 REM install IRQ handler
70 IRQ 1000
80 PRINT "TIMER RUNNING - PRESS ANY KEY"
90 GET K$ : IF K$="" THEN 90
100 POKE $BA40, 0     : REM disable timer
110 IRQ OFF
120 END
1000 REM === IRQ handler ===
1010 T = PEEK($BA41) : REM read status (clears IRQ)
1020 PRINT ".";       : REM do something each tick
1030 RETIRQ
\end{lstlisting}


% ═══════════════════════════════════════════════════════════════════════════════
\section{Music Status Registers (\$BA50--\$BA62)}
\label{sec:music-status}
% ═══════════════════════════════════════════════════════════════════════════════

The music status registers provide read-only visibility into the state of
the 14-voice music engine.  These are useful for visualizations, reactive
graphics, and checking whether a sequence has finished playing.

% --- $BA50 MusicStatus --------------------------------------------------

\mementry{BA50}{MusicStatus}{R}{0}

Playback status flags:

\bitfield{---}{---}{---}{---}{---}{---}{Mus}{SFX}

\begin{center}
\small
\begin{tabular}{@{}cl@{}}
\toprule
\textbf{Bit} & \textbf{Function} \\
\midrule
0 & SFX playing: 1 $=$ a sound effect is currently active \\
1 & Music/MIDI playing: 1 $=$ a music sequence or MIDI file is playing \\
2--7 & Reserved \\
\bottomrule
\end{tabular}
\end{center}

\begin{lstlisting}
10 REM wait for music to finish
20 IF PEEK($BA50) AND 2 THEN 20
30 PRINT "PLAYBACK COMPLETE"
\end{lstlisting}

% --- $BA51-$BA5E MusicNote -----------------------------------------------

\memrange{BA51}{BA5E}{MusicNote1--14}{R}{0}

Current MIDI note number for each of the 14 voices.  A value of 0 means
the voice is silent; values 1--127 correspond to standard MIDI note
numbers (60 $=$ Middle~C).

\begin{center}
\small
\begin{tabular}{@{}cll@{}}
\toprule
\textbf{Address} & \textbf{Voice} & \textbf{Engine} \\
\midrule
\$BA51 & Voice 0  & SID chip 1, voice 1 \\
\$BA52 & Voice 1  & SID chip 1, voice 2 \\
\$BA53 & Voice 2  & SID chip 1, voice 3 \\
\$BA54 & Voice 3  & SID chip 2, voice 1 \\
\$BA55 & Voice 4  & SID chip 2, voice 2 \\
\$BA56 & Voice 5  & SID chip 2, voice 3 \\
\$BA57 & Voice 6  & WTS voice 0 \\
\$BA58 & Voice 7  & WTS voice 1 \\
\$BA59 & Voice 8  & WTS voice 2 \\
\$BA5A & Voice 9  & WTS voice 3 \\
\$BA5B & Voice 10 & WTS voice 4 \\
\$BA5C & Voice 11 & WTS voice 5 \\
\$BA5D & Voice 12 & WTS voice 6 \\
\$BA5E & Voice 13 & WTS voice 7 \\
\bottomrule
\end{tabular}
\end{center}

\begin{lstlisting}
10 REM display active notes in real-time
20 FOR V=0 TO 13
30   N = PEEK($BA51+V)
40   IF N>0 THEN PRINT "V"; V; "="; N; " ";
50 NEXT V
60 PRINT
70 VSYNC : GOTO 20
\end{lstlisting}

% --- $BA5F-$BA60 MusicElapsed -------------------------------------------

\memrange{BA5F}{BA60}{MusicElapsed}{R}{0}

Elapsed frames since playback started, as a 16-bit little-endian value.
At 60\,Hz, divide by 60 to get elapsed seconds.

% --- $BA61-$BA62 MusicTotal ---------------------------------------------

\memrange{BA61}{BA62}{MusicTotal}{R}{0}

Total frames in the current music sequence, as a 16-bit little-endian
value.  Together with MusicElapsed, this allows programs to display a
progress indicator.

\begin{lstlisting}
10 REM display music progress
20 E = DEEK($BA5F) : T = DEEK($BA61)
30 IF T=0 THEN PRINT "NO MUSIC LOADED" : END
40 PCT = INT(E*100/T)
50 PRINT "PROGRESS:"; PCT; "%  ("; INT(E/60); "/"; INT(T/60); " SEC)"
60 VSYNC : IF PEEK($BA50) AND 2 THEN 20
70 PRINT "DONE"
\end{lstlisting}


% ═══════════════════════════════════════════════════════════════════════════════
\section{DMA Controller (\$BA63--\$BA82)}
\label{sec:dma}
% ═══════════════════════════════════════════════════════════════════════════════

The DMA Controller provides high-speed block transfers between any
combination of six memory spaces: CPU RAM, four VGC memory regions
(character, color, graphics, sprite), and XRAM.  It operates at 4 bytes
per CPU cycle, making it vastly faster than any software copy loop.

\subsection{Register Map}

\begin{center}
\small
\begin{longtable}{@{}llcll@{}}
\toprule
\textbf{Address} & \textbf{Name} & \textbf{R/W} & \textbf{Default} & \textbf{Description} \\
\midrule
\endfirsthead
\toprule
\textbf{Address} & \textbf{Name} & \textbf{R/W} & \textbf{Default} & \textbf{Description} \\
\midrule
\endhead
\$BA63 & DmaCmd       & W   & ---  & Command register \\
\$BA64 & DmaStatus    & R   & \$00 & Status code \\
\$BA65 & DmaErrCode   & R   & 0    & Error code \\
\$BA66 & DmaSrcSpace  & R/W & 0    & Source memory space \\
\$BA67 & DmaDstSpace  & R/W & 0    & Destination memory space \\
\$BA68 & DmaSrcL      & R/W & 0    & Source address low byte \\
\$BA69 & DmaSrcM      & R/W & 0    & Source address middle byte \\
\$BA6A & DmaSrcH      & R/W & 0    & Source address high byte \\
\$BA6B & DmaDstL      & R/W & 0    & Destination address low byte \\
\$BA6C & DmaDstM      & R/W & 0    & Destination address middle byte \\
\$BA6D & DmaDstH      & R/W & 0    & Destination address high byte \\
\$BA6E & DmaLenL      & R/W & 0    & Transfer length low byte \\
\$BA6F & DmaLenM      & R/W & 0    & Transfer length middle byte \\
\$BA70 & DmaLenH      & R/W & 0    & Transfer length high byte \\
\$BA71 & DmaMode      & R/W & 0    & Mode flags \\
\$BA72 & DmaFillValue & R/W & 0    & Fill byte \\
\$BA73 & DmaCountL    & R   & 0    & Bytes transferred (low) \\
\$BA74 & DmaCountM    & R   & 0    & Bytes transferred (mid) \\
\$BA75 & DmaCountH    & R   & 0    & Bytes transferred (high) \\
\bottomrule
\end{longtable}
\end{center}

\subsection{Detailed Register Descriptions}

% --- $BA63 DmaCmd -------------------------------------------------------

\mementry{BA63}{DmaCmd}{W}{---}

Command register.  Writing \$01 starts a DMA transfer using the parameters
set in the surrounding registers.

\cmdentry{01}{Start Transfer}{Begin DMA operation}

Initiates a block copy (or fill) according to the current register state.
The transfer completes before the next CPU instruction executes.

\textbf{Result:} DmaStatus $=$ \$02 on success, \$03 on error.

% --- $BA64 DmaStatus ----------------------------------------------------

\mementry{BA64}{DmaStatus}{R}{\$00}

Status of the last DMA operation:

\begin{center}
\small
\begin{tabular}{@{}cl@{}}
\toprule
\textbf{Value} & \textbf{Meaning} \\
\midrule
\$00 & Idle \\
\$01 & Busy (transfer in progress) \\
\$02 & OK --- transfer completed \\
\$03 & Error --- see DmaErrCode \\
\bottomrule
\end{tabular}
\end{center}

% --- $BA65 DmaErrCode ---------------------------------------------------

\mementry{BA65}{DmaErrCode}{R}{0}

Detailed error code when DmaStatus is \$03:

\begin{center}
\small
\begin{tabular}{@{}cl@{}}
\toprule
\textbf{Value} & \textbf{Meaning} \\
\midrule
0 & No error \\
1 & Bad command --- invalid command byte \\
2 & Bad space --- invalid source or destination space ID \\
3 & Range error --- address or length out of bounds \\
4 & Bad arguments --- invalid parameter combination \\
5 & Write protected --- destination is read-only \\
\bottomrule
\end{tabular}
\end{center}

% --- $BA66 DmaSrcSpace --------------------------------------------------

\mementry{BA66}{DmaSrcSpace}{R/W}{0}

Source memory space selector:

\begin{center}
\small
\begin{tabular}{@{}cl@{}}
\toprule
\textbf{Value} & \textbf{Space} \\
\midrule
0 & CPU RAM (64\,KB) \\
1 & VGC Character RAM (2000 bytes) \\
2 & VGC Color RAM (2000 bytes) \\
3 & VGC Graphics bitmap (32\,KB) \\
4 & VGC Sprite data \\
5 & XRAM (512\,KB, 24-bit addressing) \\
\bottomrule
\end{tabular}
\end{center}

% --- $BA67 DmaDstSpace --------------------------------------------------

\mementry{BA67}{DmaDstSpace}{R/W}{0}

Destination memory space selector.  Uses the same space IDs as DmaSrcSpace.

% --- $BA68-$BA6A DmaSrc -------------------------------------------------

\mementry{BA68}{DmaSrcL}{R/W}{0}

Low byte of the 24-bit source address.

\mementry{BA69}{DmaSrcM}{R/W}{0}

Middle byte of the source address.

\mementry{BA6A}{DmaSrcH}{R/W}{0}

High byte of the source address.  For CPU RAM (space 0), only the low two
bytes are significant.  For XRAM (space 5), all three bytes are used.

% --- $BA6B-$BA6D DmaDst -------------------------------------------------

\mementry{BA6B}{DmaDstL}{R/W}{0}

Low byte of the 24-bit destination address.

\mementry{BA6C}{DmaDstM}{R/W}{0}

Middle byte of the destination address.

\mementry{BA6D}{DmaDstH}{R/W}{0}

High byte of the destination address.

% --- $BA6E-$BA70 DmaLen -------------------------------------------------

\mementry{BA6E}{DmaLenL}{R/W}{0}

Low byte of the 24-bit transfer length.

\mementry{BA6F}{DmaLenM}{R/W}{0}

Middle byte of the transfer length.

\mementry{BA70}{DmaLenH}{R/W}{0}

High byte of the transfer length.

% --- $BA71 DmaMode ------------------------------------------------------

\mementry{BA71}{DmaMode}{R/W}{0}

DMA mode flags:

\bitfield{---}{---}{---}{---}{---}{---}{---}{Fill}

\begin{center}
\small
\begin{tabular}{@{}cl@{}}
\toprule
\textbf{Bit} & \textbf{Function} \\
\midrule
0 & Fill mode: when set, the destination is filled with DmaFillValue
    instead of copying from the source \\
1--7 & Reserved (write 0) \\
\bottomrule
\end{tabular}
\end{center}

% --- $BA72 DmaFillValue -------------------------------------------------

\mementry{BA72}{DmaFillValue}{R/W}{0}

The byte value used to fill the destination when DmaMode bit~0 is set.

% --- $BA73-$BA75 DmaCount -----------------------------------------------

\mementry{BA73}{DmaCountL}{R}{0}

Low byte of the 24-bit count of bytes transferred so far.

\mementry{BA74}{DmaCountM}{R}{0}

Middle byte of the bytes transferred count.

\mementry{BA75}{DmaCountH}{R}{0}

High byte of the bytes transferred count.

\subsection{BASIC Examples}

\subsubsection{Copy Screen RAM to XRAM}

\begin{lstlisting}
10 REM save a copy of screen RAM to XRAM bank 0
20 POKE $BA66, 1      : REM source = VGC char RAM
30 POKE $BA67, 5      : REM destination = XRAM
40 REM source offset = 0
50 POKE $BA68, 0 : POKE $BA69, 0 : POKE $BA6A, 0
60 REM destination = XRAM $000000
70 POKE $BA6B, 0 : POKE $BA6C, 0 : POKE $BA6D, 0
80 REM length = 2000 bytes (80x25 screen)
90 POKE $BA6E, $D0 : POKE $BA6F, $07 : POKE $BA70, 0
100 POKE $BA71, 0     : REM copy mode (not fill)
110 POKE $BA63, $01   : REM start transfer
120 IF PEEK($BA64)=$02 THEN PRINT "SCREEN SAVED TO XRAM"
130 IF PEEK($BA64)=$03 THEN PRINT "ERROR:"; PEEK($BA65)
\end{lstlisting}

\subsubsection{Fill Graphics Bitmap with a Color}

\begin{lstlisting}
10 REM fill the entire graphics bitmap with color $05
20 POKE $BA67, 3      : REM destination = VGC graphics
30 REM destination offset = 0
40 POKE $BA6B, 0 : POKE $BA6C, 0 : POKE $BA6D, 0
50 REM length = 32000 bytes (320x200 / 2 nybbles)
60 POKE $BA6E, 0 : POKE $BA6F, $7D : POKE $BA70, 0
70 POKE $BA71, 1      : REM fill mode
80 POKE $BA72, $55    : REM fill value ($5 in both nybbles)
90 POKE $BA63, $01    : REM start transfer
100 IF PEEK($BA64)=$02 THEN PRINT "GRAPHICS CLEARED"
\end{lstlisting}

\begin{tipbox}
The BASIC statements \texttt{DMACOPY} and \texttt{DMAFILL} provide
convenient wrappers around the DMA controller.  Direct register access
is useful when you need 24-bit XRAM addressing or fine-grained error
checking.
\end{tipbox}


% ═══════════════════════════════════════════════════════════════════════════════
\section{Blitter Controller (\$BA83--\$BAA2)}
\label{sec:blitter}
% ═══════════════════════════════════════════════════════════════════════════════

The Blitter Controller performs 2D rectangular block transfers with stride
control.  Unlike the DMA controller which operates on flat linear regions,
the blitter understands rows and columns, making it ideal for copying
rectangular regions of screen or bitmap data.  It supports fill mode and
color-key transparency for sprite-like compositing.

The blitter operates at 2 operations per CPU cycle.

\subsection{Register Map}

\begin{center}
\small
\begin{longtable}{@{}llcll@{}}
\toprule
\textbf{Address} & \textbf{Name} & \textbf{R/W} & \textbf{Default} & \textbf{Description} \\
\midrule
\endfirsthead
\toprule
\textbf{Address} & \textbf{Name} & \textbf{R/W} & \textbf{Default} & \textbf{Description} \\
\midrule
\endhead
\$BA83 & BltCmd       & W   & ---  & Command register \\
\$BA84 & BltStatus    & R   & \$00 & Status code \\
\$BA85 & BltErrCode   & R   & 0    & Error code \\
\$BA86 & BltSrcSpace  & R/W & 0    & Source memory space \\
\$BA87 & BltDstSpace  & R/W & 0    & Destination memory space \\
\$BA88 & BltSrcL      & R/W & 0    & Source base address low \\
\$BA89 & BltSrcM      & R/W & 0    & Source base address middle \\
\$BA8A & BltSrcH      & R/W & 0    & Source base address high \\
\$BA8B & BltDstL      & R/W & 0    & Destination base address low \\
\$BA8C & BltDstM      & R/W & 0    & Destination base address middle \\
\$BA8D & BltDstH      & R/W & 0    & Destination base address high \\
\$BA8E & BltWidthL    & R/W & 0    & Rectangle width low \\
\$BA8F & BltWidthH    & R/W & 0    & Rectangle width high \\
\$BA90 & BltHeightL   & R/W & 0    & Rectangle height low \\
\$BA91 & BltHeightH   & R/W & 0    & Rectangle height high \\
\$BA92 & BltSrcStrideL & R/W & 0   & Source row stride low \\
\$BA93 & BltSrcStrideH & R/W & 0   & Source row stride high \\
\$BA94 & BltDstStrideL & R/W & 0   & Destination row stride low \\
\$BA95 & BltDstStrideH & R/W & 0   & Destination row stride high \\
\$BA96 & BltMode      & R/W & 0    & Mode flags \\
\$BA97 & BltFillValue & R/W & 0    & Fill byte \\
\$BA98 & BltColorKey  & R/W & 0    & Transparent color \\
\$BA99 & BltCountL    & R   & 0    & Bytes written (low) \\
\$BA9A & BltCountM    & R   & 0    & Bytes written (middle) \\
\$BA9B & BltCountH    & R   & 0    & Bytes written (high) \\
\bottomrule
\end{longtable}
\end{center}

\subsection{Detailed Register Descriptions}

% --- $BA83 BltCmd -------------------------------------------------------

\mementry{BA83}{BltCmd}{W}{---}

Command register.  Writing \$01 starts a blit operation.

\cmdentry{01}{Start}{Begin blit operation}

Initiates a 2D block transfer (or fill) using the current register state.
For each row, the blitter copies (or fills) BltWidth bytes from the source
to the destination, then advances both pointers by their respective strides.
This repeats for BltHeight rows.

\textbf{Result:} BltStatus $=$ \$02 on success, \$03 on error.

% --- $BA84 BltStatus ----------------------------------------------------

\mementry{BA84}{BltStatus}{R}{\$00}

Status of the last blit operation.  Uses the same codes as the DMA
controller:

\begin{center}
\small
\begin{tabular}{@{}cl@{}}
\toprule
\textbf{Value} & \textbf{Meaning} \\
\midrule
\$00 & Idle \\
\$01 & Busy \\
\$02 & OK \\
\$03 & Error \\
\bottomrule
\end{tabular}
\end{center}

% --- $BA85 BltErrCode ---------------------------------------------------

\mementry{BA85}{BltErrCode}{R}{0}

Error code.  Uses the same codes as the DMA controller (see
Section~\ref{sec:dma}).

% --- $BA86-$BA87 BltSrcSpace/BltDstSpace --------------------------------

\mementry{BA86}{BltSrcSpace}{R/W}{0}

Source memory space selector.  Uses the same space IDs as the DMA controller
(0 $=$ CPU RAM, 1 $=$ VGC char, 2 $=$ VGC color, 3 $=$ VGC gfx, 4 $=$ VGC
sprite, 5 $=$ XRAM).

\mementry{BA87}{BltDstSpace}{R/W}{0}

Destination memory space selector.

% --- $BA88-$BA8A BltSrc -------------------------------------------------

\mementry{BA88}{BltSrcL}{R/W}{0}

Low byte of the 24-bit source base address.

\mementry{BA89}{BltSrcM}{R/W}{0}

Middle byte of the source base address.

\mementry{BA8A}{BltSrcH}{R/W}{0}

High byte of the source base address.

% --- $BA8B-$BA8D BltDst -------------------------------------------------

\mementry{BA8B}{BltDstL}{R/W}{0}

Low byte of the 24-bit destination base address.

\mementry{BA8C}{BltDstM}{R/W}{0}

Middle byte of the destination base address.

\mementry{BA8D}{BltDstH}{R/W}{0}

High byte of the destination base address.

% --- $BA8E-$BA8F BltWidth -----------------------------------------------

\mementry{BA8E}{BltWidthL}{R/W}{0}

Low byte of the rectangle width in bytes.

\mementry{BA8F}{BltWidthH}{R/W}{0}

High byte of the rectangle width.

% --- $BA90-$BA91 BltHeight ----------------------------------------------

\mementry{BA90}{BltHeightL}{R/W}{0}

Low byte of the rectangle height in rows.

\mementry{BA91}{BltHeightH}{R/W}{0}

High byte of the rectangle height.

% --- $BA92-$BA93 BltSrcStride -------------------------------------------

\mementry{BA92}{BltSrcStrideL}{R/W}{0}

Low byte of the source row stride.  The stride is the number of bytes
between the start of one row and the start of the next in the source
memory space.  For a full-width screen copy, this equals the screen width.

\mementry{BA93}{BltSrcStrideH}{R/W}{0}

High byte of the source row stride.

% --- $BA94-$BA95 BltDstStride -------------------------------------------

\mementry{BA94}{BltDstStrideL}{R/W}{0}

Low byte of the destination row stride.

\mementry{BA95}{BltDstStrideH}{R/W}{0}

High byte of the destination row stride.

% --- $BA96 BltMode ------------------------------------------------------

\mementry{BA96}{BltMode}{R/W}{0}

Blitter mode flags:

\bitfield{---}{---}{---}{---}{---}{---}{Key}{Fill}

\begin{center}
\small
\begin{tabular}{@{}cp{10cm}@{}}
\toprule
\textbf{Bit} & \textbf{Function} \\
\midrule
0 & Fill mode: when set, the destination is filled with BltFillValue
    instead of copying from the source.  Width and height still control
    the rectangular region. \\
1 & Color-key transparency: when set, source bytes matching BltColorKey
    are skipped (the destination byte is not overwritten).  This enables
    transparent blitting for sprite-like effects. \\
2--7 & Reserved (write 0) \\
\bottomrule
\end{tabular}
\end{center}

% --- $BA97 BltFillValue -------------------------------------------------

\mementry{BA97}{BltFillValue}{R/W}{0}

Fill byte used when BltMode bit~0 is set.

% --- $BA98 BltColorKey --------------------------------------------------

\mementry{BA98}{BltColorKey}{R/W}{0}

Transparent color.  When BltMode bit~1 is set, source bytes equal to this
value are skipped during the copy, leaving the destination unchanged.

% --- $BA99-$BA9B BltCount -----------------------------------------------

\mementry{BA99}{BltCountL}{R}{0}

Low byte of the 24-bit count of bytes actually written.

\mementry{BA9A}{BltCountM}{R}{0}

Middle byte.

\mementry{BA9B}{BltCountH}{R}{0}

High byte.  When color-key transparency is active, BltCount may be less
than Width $\times$ Height because skipped pixels are not counted.

\subsection{BASIC Examples}

\subsubsection{Blit a Rectangular Region with Transparency}

\begin{lstlisting}
10 REM blit a 20x10 region from XRAM to graphics
20 REM with color 0 as transparent
30 POKE $BA86, 5      : REM source = XRAM
40 POKE $BA87, 3      : REM destination = VGC graphics
50 REM source address = XRAM $000000
60 POKE $BA88, 0 : POKE $BA89, 0 : POKE $BA8A, 0
70 REM destination = graphics offset 1000
80 POKE $BA8B, $E8 : POKE $BA8C, $03 : POKE $BA8D, 0
90 REM width=20, height=10
100 DOKE $BA8E, 20 : DOKE $BA90, 10
110 REM source stride = 20 (tiles packed tightly)
120 DOKE $BA92, 20
130 REM destination stride = 160 (320/2 nybbles per row)
140 DOKE $BA94, 160
150 REM mode: color-key transparency on (bit 1)
160 POKE $BA96, 2
170 POKE $BA98, 0     : REM color key = 0 (transparent)
180 POKE $BA83, $01   : REM start blit
190 IF PEEK($BA84)=$02 THEN PRINT "BLIT DONE"
200 IF PEEK($BA84)=$03 THEN PRINT "ERROR:"; PEEK($BA85)
\end{lstlisting}

\subsubsection{Fill a Rectangle on Screen}

\begin{lstlisting}
10 REM fill a 40x5 rectangle of character RAM with 'X'
20 POKE $BA87, 1      : REM destination = VGC char RAM
30 REM destination = row 10, column 20 (offset = 10*80+20 = 820)
40 POKE $BA8B, $34 : POKE $BA8C, $03 : POKE $BA8D, 0
50 DOKE $BA8E, 40     : REM width = 40 columns
60 DOKE $BA90, 5      : REM height = 5 rows
70 DOKE $BA94, 80     : REM destination stride = 80 (screen width)
80 POKE $BA96, 1      : REM fill mode
90 POKE $BA97, ASC("X") : REM fill with 'X'
100 POKE $BA83, $01   : REM start blit
\end{lstlisting}

\begin{tipbox}
The BASIC statements \texttt{BLITCOPY} and \texttt{BLITFILL} provide
convenient wrappers around the blitter.  Use direct register access when
you need color-key transparency or cross-space blits that the BASIC
statements do not expose.
\end{tipbox}


% ═══════════════════════════════════════════════════════════════════════════════
\section{XMC Windows (\$BC00--\$BFFF)}
\label{sec:xmc-windows}
% ═══════════════════════════════════════════════════════════════════════════════

\memrange{BC00}{BFFF}{XMC Window}{R/W}{0}

The final 1\,KB of this address range is divided into four 256-byte
windows that provide direct memory-mapped access to XRAM.  Each window
maps a contiguous 256-byte page of XRAM into the 6502 address space,
allowing reads and writes to pass through transparently --- no commands,
no register setup, just ordinary \texttt{PEEK} and \texttt{POKE}.

\begin{center}
\small
\begin{tabular}{@{}llll@{}}
\toprule
\textbf{CPU Address} & \textbf{Window} & \textbf{Enable Bit} & \textbf{Base Register} \\
\midrule
\$BC00--\$BCFF & Window 0 & Bit 0 of \addr{BA16} & XmcWin0A (\addr{BA18}--\addr{BA1A}) \\
\$BD00--\$BDFF & Window 1 & Bit 1 of \addr{BA16} & XmcWin1A (\addr{BA1B}--\addr{BA1D}) \\
\$BE00--\$BEFF & Window 2 & Bit 2 of \addr{BA16} & XmcWin2A (\addr{BA1E}--\addr{BA20}) \\
\$BF00--\$BFFF & Window 3 & Bit 3 of \addr{BA16} & XmcWin3A (\addr{BA21}--\addr{BA23}) \\
\bottomrule
\end{tabular}
\end{center}

When a window is \textbf{enabled}, reads from its 256-byte range return the
corresponding byte from XRAM, and writes store the byte into XRAM.  When a
window is \textbf{disabled}, reads return \$00 and writes are silently
discarded.

The base address register specifies the XRAM address that the window's
first byte (\$xx00) maps to.  The remaining 255 bytes map to consecutive
XRAM addresses.

\begin{lstlisting}
10 REM map window 0 to XRAM page $000100
20 POKE $BA18, 0 : POKE $BA19, 1 : POKE $BA1A, 0
30 POKE $BA16, 1  : REM enable window 0
40 POKE $BC00, 42 : REM writes to XRAM $000100
50 PRINT PEEK($BC00) : REM reads from XRAM $000100
\end{lstlisting}

\begin{lstlisting}
10 REM use all four windows simultaneously
20 REM map to four consecutive XRAM pages
30 FOR W=0 TO 3
40   POKE $BA18+W*3, 0 : POKE $BA19+W*3, W : POKE $BA1A+W*3, 0
50 NEXT W
60 POKE $BA16, $0F : REM enable all four windows
70 REM now $BC00-$BFFF maps to XRAM $000000-$0003FF
80 FOR I=0 TO 1023
90   POKE $BC00+I, I AND 255
100 NEXT I
110 PRINT "WROTE 1KB TO XRAM VIA WINDOWS"
\end{lstlisting}

\begin{notebox}
Windows provide the fastest way to access small, frequently-used regions
of XRAM.  For bulk transfers, the Stash/Fetch commands or DMA controller
are more efficient.
\end{notebox}
